	\begin{center}
		\section{进阶数学}
	\end{center}

	\begin{multicols*}{2} % 开始双栏，数字 2 表示分两栏

		\subsection{FFT多项式乘法（非递归）}

		采用非递归版本实现FFT（蝴蝶变换），能显著减小常数并优化内存使用。
		这里通过 \texttt{multiply} 演示了如何运用 \texttt{fft} 对两个多项式（系数数组表示）进行乘法操作。

		\begin{minted}{cpp}
// 理论讲解 https://www.bilibili.com/video/BV1za411F76U/
// 参考（抄的） https://cp-algorithms.com/algebra/fft.html
// 多项式乘法 https://www.luogu.com.cn/record/229525769
using CD = complex<double>;

void fft(vector<CD> &a, bool invert) {
	int n = a.size();
	for (int i = 1, j = 0; i < n; i++) {
		int bit = n >> 1;
		for (; j & bit; bit >>= 1) j ^= bit;
		j ^= bit;
		if (i < j) swap(a[i], a[j]);
	}
	for (int len = 2; len <= n; len <<= 1) {
		double ang = 2 * pi / len * (invert ? -1 : 1);
		CD wlen(cos(ang), sin(ang));
		for (int i = 0; i < n; i += len) {
			CD w(1);
			for (int j = 0; j < len / 2; j++) {
				CD u = a[i + j], v = a[i + j + len / 2] * w;
				a[i + j] = u + v;
				a[i + j + len / 2] = u - v;
				w *= wlen;
			}
		}
	}
	if (invert) {
		for (CD &x : a) x /= n;
	}
}

vector<ll> multiply(vector<ll> &a, vector<ll> &b) {
	vector<CD> fa(a.begin(), a.end()), fb(b.begin(), b.end());
	int n = 1;
	while (n < (int)(a.size() + b.size())) {
		n <<= 1;
	}
	fa.resize(n); fb.resize(n);
	fft(fa, false); fft(fb, false);
	for (int i = 0; i < n; i++) fa[i] *= fb[i];
	fft(fa, true);
	vector<ll> res(n);
	// 避免精度损失
	for (int i = 0; i < n; i++) res[i] = llround(fa[i].real());
	return res;
}
		\end{minted}

		\subsubsection{高精度乘法（基于FFT）}

		https://www.luogu.com.cn/record/229526218

		大整数乘法的核心思想：将大整数的每一位看作多项式的一个系数。例如 $1234$ 可以表示为多项式 $4 + 3x + 2x^2 + x^3$（低位在前）。两个大整数相乘等价于两个多项式相乘，乘完后对结果数组统一处理进位。

		\begin{minted}{cpp}
inline void __() {
	string A, B;
	cin >> A >> B;
	vector<ll> a, b;
	for (int i = (int)A.size() - 1; i >= 0; --i) {
		a.push_back(A[i] - '0');
	}
	for (int i = (int)B.size() - 1; i >= 0; --i) {
		b.push_back(B[i] - '0');
	}
	
	vector<ll> ans = multiply(a, b);
	ll carry = 0;
	for (int i = 0; i < (int)ans.size(); ++i) {
		ans[i] += carry;
		carry = ans[i] / 10;
		ans[i] %= 10;
	}
	while (!ans.empty() && ans.back() == 0) {
		ans.pop_back();
	}
	if (ans.empty()) {
		cout << 0;
		return;
	}
	for (int i = (int)ans.size() - 1; i >= 0; --i) {
		cout << ans[i];
	}
}
		\end{minted}

		\subsubsection{FFT 的使用方法与常见应用}

		\texttt{multiply(a, b)} 的本质是：给定多项式 $A(x) = \sum a_i x^i$ 和 $B(x) = \sum b_i x^i$，返回乘积多项式 $C(x) = A(x) \cdot B(x)$ 的系数数组，其中 $c_k = \sum_{i+j=k} a_i \cdot b_j$。因此，\textbf{只要问题能转化为这种卷积形式，就能用 FFT 加速到 $O(n \log n)$}。

		\textbf{1. 加法卷积计数}

		最典型的场景：有两组数，问它们各取一个数相加能得到每个和的方案数。

		例如掷两颗骰子，求每种点数之和的方案数。将骰子 A 的各面出现次数存入 \texttt{a[]}（\texttt{a[k]} 表示点数 $k$ 出现的次数），骰子 B 同理存入 \texttt{b[]}，则 \texttt{multiply(a, b)} 的结果 \texttt{c[k]} 就是和为 $k$ 的方案数。

		\begin{minted}{cpp}
// 例：两个数组各取一个数，统计每种和的方案数
// a[i] = 数 i 在第一组中出现的次数
// b[j] = 数 j 在第二组中出现的次数
// c[k] = 和为 k 的方案数 = sum(a[i]*b[k-i])
vector<ll> c = multiply(a, b);
		\end{minted}

		\textbf{2. 字符串匹配（含通配符）}

		给定文本串 $T$（长度 $n$）和模式串 $P$（长度 $m$），对于 $T$ 的每个起始位置 $t$，定义失配度：
		$$D(t) = \sum_{j=0}^{m-1} (T_{t+j} - P_j)^2 \cdot w_j$$
		其中 $w_j$ 是通配符权重（通配符处 $w_j = 0$）。$D(t) = 0$ 表示位置 $t$ 完全匹配。

		将 $P$ 翻转为 $P'_j = P_{m-1-j}$，展开后各项都可化为卷积形式，分别用 FFT 计算后合并。

		\begin{minted}{cpp}
// 无通配符的简单情形：检查 T 和 P 是否在每个位置匹配
// 构造 a[i] = T[i], b[j] = P[m-1-j]（翻转 P）
// 则 c[t+m-1] = sum(T[t+j] * P[m-1-j]) 即为位置 t 的内积
// 更一般地，含通配符时需要构造多组卷积分别计算
vector<ll> a(n), b(m);
for (int i = 0; i < n; i++) a[i] = T[i] - 'a' + 1;
for (int j = 0; j < m; j++) b[j] = P[m - 1 - j] - 'a' + 1;
vector<ll> c = multiply(a, b);
// c[t + m - 1] 即为位置 t 的匹配内积
		\end{minted}

		\textbf{3. 高精度乘法}

		如上一节所述，将大整数每一位作为多项式系数，调用 \texttt{multiply} 后统一处理进位。

		\textbf{4. 多项式求值与插值的中间步骤}

		分治 FFT 可以加速多项式多点求值、多点插值等操作，核心都是递归地调用 \texttt{multiply} 合并子问题的结果多项式。

	\end{multicols*} % 线性规划一节改用单栏：字典表代码行较长，双栏会频繁折行

		\subsection{线性规划（两阶段单纯形 Simplex）}

		求解标准型
		$$
		\max\ c^{\mathsf T} z \quad \text{s.t.} \quad A z \le b,\ z \ge 0 .
		$$
		三样输入分别是：$A$ 是\textbf{约束系数}（$A_{ij}$ = 第 $i$ 条约束里变量 $j$ 的系数），$b$ 是每条约束的\textbf{上限}，$c$ 是\textbf{要最大化的那个式子的系数}。代码里就直接叫 \texttt{coef} / \texttt{limit} / \texttt{profit}，解向量叫 \texttt{x}。

		可行域是若干半空间交出的凸多面体，最优值必在某个\textbf{顶点}处取到。单纯形法就是从一个顶点出发，每次沿一条棱走到相邻的更优顶点（一次 \texttt{pivot} = 一个非基变量入基、一个基变量出基），直到目标行里没有改进方向为止。

		\begin{center}
		\begin{tikzpicture}[scale=1.5, every node/.style={font=\small}]
			\coordinate (O) at (0,0);
			\coordinate (A) at (3,0);
			\coordinate (B) at (3,1);
			\coordinate (C) at (1.5,2.5);
			\coordinate (D) at (0,3);
			% 约束直线（延伸到可行域外用浅色虚线）
			\draw[gray!55, dashed] (0,4) -- (4,0);
			\draw[gray!55, dashed] (0,3) -- (4,1.667);
			\draw[gray!55, dashed] (3,-0.25) -- (3,3.3);
			\fill[blue!9] (O) -- (A) -- (B) -- (C) -- (D) -- cycle;
			\draw[blue!60, thick] (O) -- (A) -- (B) -- (C) -- (D) -- cycle;
			\draw[->] (-0.25,0) -- (4.5,0) node[below] {$x$};
			\draw[->] (0,-0.25) -- (0,3.8) node[left] {$y$};
			% 约束名
			\node[font=\scriptsize, gray!70!black] at (4.05,0.35) {$x+y\le 4$};
			\node[font=\scriptsize, gray!70!black] at (4.05,1.95) {$x+3y\le 9$};
			\node[font=\scriptsize, gray!70!black] at (3.0,3.6) {$x\le 3$};
			% 单纯形游走路径
			\draw[->, red!80!black, very thick] (O) -- (A);
			\draw[->, red!80!black, very thick] (A) -- (B);
			\foreach \p in {O,A,B,C,D} \fill (\p) circle (1.7pt);
			\node[font=\scriptsize, below left=1pt and 1pt] at (O) {$z=0$};
			\node[font=\scriptsize, below=5pt] at (A) {$z=9$};
			\node[font=\scriptsize, right=4pt, red!70!black] at (B) {$z=11$};
			\node[font=\scriptsize, above right=1pt and 2pt] at (C) {$z=9.5$};
			\node[font=\scriptsize, above left=1pt and 2pt] at (D) {$z=6$};
		\end{tikzpicture}
		\end{center}

		图中目标为 $\max z = 3x + 2y$，路径 $(0,0) \to (3,0) \to (3,1)$ 即两次 pivot 后到达最优顶点。

		\textbf{实现要点}：

		\begin{itemize}
			\item \textbf{字典表连续存储}：$(rows+2) \times (cols+2)$ 一维数组行主序，比 \texttt{vector<vector<double>>} 少一层间接寻址，内层循环无分支可向量化。
			\item \textbf{两阶段}：$b$ 有负分量时原点不可行，先引入人工变量 $x_0$ 跑第一阶段 $\min x_0$；若最优 $x_0 > 0$ 则原问题\textbf{不可行}。$b \ge 0$ 时直接跳过第一阶段。
			\item \textbf{Devex 定价}：入基列挑 $d_j^2 / w_j$ 最大者（$d_j$ 为目标行系数、$w_j$ 为参考权重），比朴素 Dantzig 规则（挑最负的 $d_j$）迭代次数少很多。
			\item \textbf{防循环}：退化时目标值长期不推进，连续 \texttt{stallLimit} 次后切 Bland 最小下标规则（保证有限步终止），推进后自动切回 Devex。
			\item \textbf{出基行}：最小比值 $b_i / a_{is}$；比值持平时选主元绝对值更大的一行，数值更稳。
		\end{itemize}

		\textbf{返回值}：最优值；无界返回 $+\infty$，不可行返回 $-\infty$；\texttt{solution} 非空时回填一组最优解。

		\textbf{建模转换}：等式约束拆成 $\le$ 与 $\ge$ 两条；$\ge$ 约束两边取负变 $\le$；$\min$ 问题把 \texttt{profit} 整体取负、答案再取负；变量可正可负时拆成 $z = z^{+} - z^{-}$ 两个非负变量。

		\begin{minted}{cpp}
// 两个变量 x、y，都 >= 0；要最大化 3x + 2y
// 三条约束：x + y <= 4，x + 3y <= 9，x <= 3
vector<vector<double>> coef = {{1, 1},      // 第 1 条约束里 x、y 的系数
                               {1, 3},      // 第 2 条
                               {1, 0}};     // 第 3 条（y 不出现，系数写 0）
vector<double> limit = {4, 9, 3};           // 每条约束不等号右边的上限
vector<double> profit = {3, 2};             // 目标 3x + 2y 里 x、y 的系数
SimplexLP lp;
lp.init(coef, limit, profit);
vector<double> x;                           // 取到最优时每个变量的值
double best = lp.solve(&x);					// best = 11, x = {3, 1}
		\end{minted}

		\begin{minted}{cpp}
// 变量 x[0..cols) 全都 >= 0；每条约束形如 coef[i][0]*x[0] + ... <= limit[i]；
// 目标是让 profit[0]*x[0] + profit[1]*x[1] + ... 尽量大。
struct SimplexLP {
    // 核心就是一张表 tab，(rows + 2) 行 x (cols + 2) 列，拍平成一维存（比二维 vector 少一层寻址）：
    //   每行一条约束，每列一个变量，最后一列放这条约束还剩多少余量
    //   多出的倒数第 2 行 = 目标（每个变量再加一单位能赚多少），最后一行只在"先找可行点"时用
    //   多出的倒数第 2 列 = 辅助变量 x0，同样只在"先找可行点"时用
    int rows = 0;               // 约束条数
    int cols = 0;               // 变量个数
    int stride = 0;             // 一行占多少个 double（= cols + 2）
    vector<double> tab;
    vector<int> basis;          // basis[i] = 第 i 行当前解出来的那个变量的编号
    vector<int> nonbasis;       // nonbasis[j] = 第 j 列对应的变量编号，这些变量当前取 0
    double eps = 1e-9;
    long long pivots = 0;       // 换基次数，调试看规模用
    vector<double> weight;      // Devex 参考权重，挑变量时用

    double* row(int i) { return tab.data() + static_cast<size_t>(i) * stride; }
    const double* row(int i) const { return tab.data() + static_cast<size_t>(i) * stride; }

    // coef[i][j]：第 i 条约束里变量 j 的系数     limit[i]：第 i 条约束的上限
    // profit[j] ：变量 j 每多一单位，目标增加多少
    void init(const vector<vector<double>>& coef,
              const vector<double>& limit,
              const vector<double>& profit) {
        rows = static_cast<int>(limit.size());
        cols = static_cast<int>(profit.size());
        stride = cols + 2;
        tab.assign(static_cast<size_t>(rows + 2) * stride, 0.0);
        basis.resize(rows);
        nonbasis.resize(cols + 1);
        for (int i = 0; i < rows; ++i) {
            double* p = row(i);
            for (int j = 0; j < cols; ++j) p[j] = coef[i][j];
            p[cols] = -1.0;          // 辅助变量 x0 那一列，只在找可行点时起作用
            p[cols + 1] = limit[i];
            basis[i] = cols + i;     // 一开始所有变量取 0，这一行剩的余量就是 limit[i]
        }
        double* obj = row(rows);
        // 目标行存的是 -profit：以后看到负数就说明"这个变量还能加，目标还能变大"
        for (int j = 0; j < cols; ++j) { obj[j] = -profit[j]; nonbasis[j] = j; }
        nonbasis[cols] = -1;
        row(rows + 1)[cols] = 1.0;   // 找可行点那一阶段的目标：把 x0 压到 0
        weight.assign(cols + 1, 1.0);
        pivots = 0;
    }

    // 让第 s 列的变量从 0 开始往大加，加到第 r 条约束把它卡住为止；这两个变量交换身份
    void pivot(int r, int s) {
        ++pivots;
        double* pr = row(r);
        const double inv = 1.0 / pr[s];
        {   // Devex：顺手更新参考权重，供下次挑变量用（O(cols)）
            const double as = pr[s], ws = weight[s];
            for (int j = 0; j <= cols; ++j) {
                if (j == s) continue;
                const double ratio = pr[j] / as;
                const double cand = ratio * ratio * ws;
                if (cand > weight[j]) weight[j] = cand;
            }
            const double nw = ws / (as * as);
            weight[s] = nw > 1.0 ? nw : 1.0;
        }
        const int total = rows + 2, width = cols + 2;
        for (int i = 0; i < total; ++i) {   // 用第 r 行把其余每行的第 s 列消成 0
            if (i == r) continue;
            double* pi = row(i);
            const double f = pi[s] * inv;
            if (f == 0.0) continue;
            for (int j = 0; j < width; ++j) pi[j] -= pr[j] * f;   // 内层无分支，编译器可向量化
            pi[s] = -f;                                           // 第 s 列要单独修正
        }
        for (int j = 0; j < width; ++j) pr[j] *= inv;
        pr[s] = inv;
        swap(basis[r], nonbasis[s]);
    }

    // phase = 1：先找一个满足所有约束的点（把 x0 压到 0）
    // phase = 2：在可行的范围内把目标做到最大
    bool run(int phase) {
        const int objRow = (phase == 1) ? rows + 1 : rows;
        // 目标值长时间不动（退化）就换 Bland 规则：它慢，但保证不会绕圈死循环
        long long stall = 0;
        const long long stallLimit = 4LL * (rows + cols) + 64;
        while (true) {
            const double* obj = row(objRow);
            int s = -1;                                  // 先挑"加谁"：哪个变量还能让目标变大
            if (stall < stallLimit) {                    // Devex：挑性价比最高的
                double best = 0.0;
                for (int j = 0; j <= cols; ++j) {
                    if (phase == 2 && nonbasis[j] == -1) continue;
                    const double d = obj[j];
                    if (d > -eps) continue;              // 不是负的 => 加它目标也不会变大
                    const double score = d * d / weight[j];
                    if (score > best) { best = score; s = j; }
                }
            } else {                                     // Bland：老实挑编号最小的
                for (int j = 0; j <= cols; ++j) {
                    if (phase == 2 && nonbasis[j] == -1) continue;
                    if (obj[j] < -eps && (s == -1 || nonbasis[j] < nonbasis[s])) s = j;
                }
            }
            if (s == -1) return true;                    // 没变量能让目标变大 => 这一阶段做完

            // 再看"加到哪停"：哪条约束最先被顶满，即余量 / 系数 最小的那一行
            int r = -1;
            double bestRatio = 0.0, bestPiv = 0.0;
            for (int i = 0; i < rows; ++i) {
                const double* pi = row(i);
                if (pi[s] < eps) continue;               // 系数 <= 0，这条约束拦不住它
                const double ratio = pi[cols + 1] / pi[s];
                if (r == -1 || ratio < bestRatio - 1e-12 ||
                    // 两行一样紧时选系数绝对值更大的那行，算起来数值更稳
                    (ratio < bestRatio + 1e-12 && pi[s] > bestPiv)) {
                    r = i; bestRatio = ratio; bestPiv = pi[s];
                }
            }
            if (r == -1) return false;                   // 没人拦得住它 => 目标能到无穷大

            const double before = row(objRow)[cols + 1];
            pivot(r, s);
            const double after = row(objRow)[cols + 1];
            stall = (after > before + 1e-12) ? 0 : stall + 1;
        }
    }

    // 返回目标能达到的最大值；无界返回 +inf，一个可行点都没有返回 -inf
    // 传了 solution 就顺便把最优时每个变量的取值填回去
    double solve(vector<double>* solution = nullptr) {
        int r = 0;
        for (int i = 1; i < rows; ++i)
            if (row(i)[cols + 1] < row(r)[cols + 1]) r = i;
        if (row(r)[cols + 1] < -eps) {                   // 有约束上限是负的 => 全取 0 都不合法，先找可行点
            pivot(r, cols);
            if (!run(1) || row(rows + 1)[cols + 1] < -eps)
                return -numeric_limits<double>::infinity();   // x0 压不到 0 => 无解
            for (int i = 0; i < rows; ++i) {
                if (basis[i] != -1) continue;            // 把还赖在解里的 x0 换出去
                int s = -1;
                for (int j = 0; j <= cols; ++j)
                    if (s == -1 || row(i)[j] < row(i)[s] ||
                        (row(i)[j] == row(i)[s] && nonbasis[j] < nonbasis[s])) s = j;
                pivot(i, s);
            }
        }
        if (!run(2)) return numeric_limits<double>::infinity();
        if (solution) {                                  // 没被解出来的变量取值就是 0
            solution->assign(cols, 0.0);
            for (int i = 0; i < rows; ++i)
                if (basis[i] < cols) (*solution)[basis[i]] = row(i)[cols + 1];
        }
        return row(rows)[cols + 1];
    }
};
		\end{minted}

